%% ---------------------------------------------------------------------
%%  FIGURA - Laco de animacao com acumulador (DG-05)
%% ---------------------------------------------------------------------
\begin{figure}[htb]
	\caption{Laço de animação com acumulador de tempo e execução de passos}
	\label{fig:game-loop}
	\centering
	\begin{tikzpicture}[
			font=\footnotesize,
			passo/.style={draw, rounded corners=2pt, minimum width=3.5cm, minimum height=0.8cm, align=center, fill=black!4},
			dec/.style={draw, diamond, aspect=2.4, inner sep=1pt, align=center, font=\scriptsize},
			>=Stealth,
			fl/.style={->, thick},
			rot/.style={font=\scriptsize, inner sep=2pt}
		]

		\node[passo] (raf) at (0,0) {Requisição de quadro};
		\node[passo, below=0.55cm of raf] (delta) {Calcular intervalo $\Delta$ e limitar a 250\,ms};
		\node[dec, below=0.75cm of delta] (rodando) {estado =\\ \texttt{running}?};
		\node[passo, right=2.5cm of rodando] (acum) {Acumular: $a \mathrel{+}= \Delta$};
		\node[dec, below=0.85cm of acum] (suf) {$a \geq t$ e\\ passos $< 5$?};
		\node[passo, right=2.3cm of suf] (tick) {Executar passo\\ de simulação};
		\node[passo, below=0.7cm of tick] (dec2) {Subtrair $t$ de $a$\\ e incrementar contador};
		\node[passo, below=2.2cm of rodando] (desc) {Descartar acúmulo residual};
		\node[passo, below=0.7cm of desc] (render) {Renderizar};
		\node[passo, below=0.7cm of render] (loop) {Agendar próximo quadro};

		\draw[fl] (raf) -- (delta);
		\draw[fl] (delta) -- (rodando);
		\draw[fl] (rodando) -- node[rot, above] {sim} (acum);
		\draw[fl] (acum) -- (suf);
		\draw[fl] (suf) -- node[rot, above] {sim} (tick);
		\draw[fl] (tick) -- (dec2);
		\draw[fl] (dec2) |- (acum);
		\draw[fl] (suf) |- node[rot, pos=0.75, left] {não} (desc);
		\draw[fl] (rodando) -- node[rot, left] {não} (desc);
		\draw[fl] (desc) -- (render);
		\draw[fl] (render) -- (loop);
		\draw[fl] (loop) -- ++(2.2,0) |- (raf);
	\end{tikzpicture}
	\legend{Fonte: elaborado pelo autor (2026) a partir de \texttt{index.html}, seção~10.
	O teto de cinco passos por quadro e o descarte do acúmulo residual são as duas defesas contra a
	espiral da morte.}
\end{figure}
